\documentclass[10pt,a4paper]{amsart}
\usepackage[margin=19mm]{geometry}
\usepackage{amsmath,amssymb,mathtools}
\usepackage[scr=rsfs,cal=euler]{mathalfa}
\usepackage{xfrac}
\usepackage[unicode,hidelinks,bookmarks=false]{hyperref}
\newcommand{\ie}{i.e.\ }
\newcommand{\cf}{cf.\ }
\newcommand{\resp}{respectively\ }
\newcommand{\Subsection}{\subsection}
\providecommand{\nobreakdash}{\nobreak\hbox{-}}
\setlength{\emergencystretch}{2em}
\multlinegap0pt
\usepackage{xcolor}
\usepackage{tikz}
\usepackage{amssymb}
\usepackage{mathtools}
\usepackage[shortlabels]{enumitem}
\usepackage{bm}
\newtheorem{lemma}{Lemma}
\usepackage{cleveref}
\crefname{lemma}{Lemma}{Lemmas}
\newcommand{\Z}{\mathbb{Z}}
\title{Khovanov homology \break and refined bounds for Gordian distances}
\author{Lukas Lewark, Laura Marino, Claudius Zibrowius}
\date{Source: arXiv:2409.05743v2 (CC BY 4.0)}
\begin{document}
\maketitle

\noindent\emph{Source and licence.} Khovanov homology \break and refined bounds for Gordian distances. Version arXiv:2409.05743v2 (CC BY 4.0), distributed by arXiv under the Creative Commons Attribution 4.0 International licence, http://creativecommons.org/licenses/by/4.0/, as recorded in the arXiv licence metadata for this exact version and shown on the abstract page https://arxiv.org/abs/2409.05743v2. Full licence notice: \texttt{licenses/arxiv-2409.05743-CC-BY-4.0.txt}.\par\smallskip

\noindent\textit{Editorial scope.} Counted unit: the article's lemma lem:optimal\_torus\_PRTR (source0.tex lines 3023--3025). The article's complete original proof of it is delivered with it (source0.tex lines 3027--3090, one \texttt{proof} environment of 64 lines). No mathematics is written, repaired or re-derived: the statement and the proof are the article's own lines, in the article's own notation, and no retained line is substituted or re-flowed.  No further source-local context is needed: every cross-reference of every retained line resolves inside the two delivered spans.  Nothing was omitted from any retained span, so the package carries no omission record.  No retained line contains a citation, so the wrapper carries no bibliography and no bibliography entry.  No retained line contains a figure environment, an image-inclusion command or drawing code, so no image is delivered and the package declares no asset dependency.  Editorial formatting only, in addition: an AMS \texttt{amsart} preamble that restates the article's own declarations and macros used by the retained lines, namely the environments \texttt{lemma} on separate counters, together with the standard packages those lines need.  The retained environments print in this document as Lemma 1 in order of appearance, whereas the article prints the same items with the numbering of its own full document.  The complete native source of the article as distributed in the arXiv e-print for this exact version is bundled unchanged as \texttt{sources/165001/source0.tex}.
\begin{lemma}\label{lem:optimal_torus_PRTR}
	For any \(n\in\Z_{\geq 0}\) there exists a knot \(K_n\) that admits a diagram \(D_n\) with \(c_+(D_n)=4n+1\) and that is related to both \(T(3,3n+1)\) and \(T(3,3n+2)\) by a proper rational tangle replacement.
\end{lemma}

\begin{proof}
    Let $B_3=\langle \sigma_1,\sigma_2 \;|\; \sigma_1\sigma_2\sigma_1 = \sigma_2\sigma_1\sigma_2 \rangle$ be the braid group on three strands.
    Let $D_n$ be the knot diagram given by the closure of the braid $\beta_n=\sigma_1^{-1}(\sigma_2\sigma_1^2\sigma_2)^n \sigma_2$. The diagram $D_n$ has $4n+1$ positive crossings.
    
    We now construct a proper rational tangle replacement to $T(3,3n+1)$. For this we need the following identity: 
    \begin{equation}\label{eq:2n+1_twists}
        (\sigma_1 \sigma_2)^{3n+1}=\sigma_1^{2n+2}\beta_n.
    \end{equation}
    We prove this identity by induction on $n\in\Z_{\geq 0}$. 
    For $n=0$ the claim is obvious.
    Now suppose \cref{eq:2n+1_twists} holds for $n$; we prove it for $n+1$. 
    Let $\Delta^2$ denote the full twist $\Delta^2=(\sigma_1\sigma_2)^3$. 
    Then 
    \[
    (\sigma_1 \sigma_2)^{3n+4}
    =
    \Delta^2(\sigma_1 \sigma_2)^{3n+1}
    =
    \Delta^2\sigma_1^{2n+2}\beta_n
    =
    \sigma_1^{2n+2}\Delta^2\beta_n.
    \]
    The last equality comes from the well-known fact that the full twist $\Delta^2$ is in the centre of the braid group.
    Moreover, observe that 
    \[
    \Delta^2
    =
    \sigma_1\sigma_2\sigma_1(\sigma_2\sigma_1\sigma_2
)    =
    \sigma_1\sigma_2\sigma_1(\sigma_1\sigma_2\sigma_1)    
    =
    \sigma_1(\sigma_2\sigma_1^2\sigma_2)\sigma_1.
    \]
    Hence
    \[
    (\sigma_1 \sigma_2)^{3n+4}
    =
    \sigma_1^{2n+2}\Delta^2\beta_n
    =
    \sigma_1^{2n+3}(\sigma_2\sigma_1^2\sigma_2)\sigma_1\beta_n
    =
    \sigma_1^{2n+4}\beta_{n+1}.
    \]

    The knot $T(3,3n+1)$ is given by the closure of the braid $(\sigma_1 \sigma_2)^{3n+1}\in B_3$.
    Observe that the operation $(\sigma_1 \sigma_2)^{3n+1}=\sigma_1^{2n+2}\beta_n \leadsto \beta_n$ corresponds to a proper rational tangle replacement $Q_{-2n-2} \leadsto Q_{0}$ from a diagram of $T(3,3n+1)$ to $D_n$. 

    A proper rational replacement for $T(3,3n+2)$ can be constructed similarly.
		We have
    \[
    (\sigma_1 \sigma_2)^{3n+2}
    =
    \sigma_1(\sigma_2 \sigma_1)^{3n} \sigma_2\sigma_1\sigma_2
    =
    \sigma_1(\Delta^2)^{n} \sigma_1\sigma_2\sigma_1
    =
    \sigma_1(\sigma_1 \sigma_2)^{3n+1}\sigma_1
    =
    \sigma_1^{2n+3}\beta_n\sigma_1.
    \]
    The knot $T(3,3n+2)$ is given by the closure of the braid $(\sigma_1 \sigma_2)^{3n+2}\in B_3$.
    Furthermore, the closures of $\sigma_1^{2n+3}\beta_n\sigma_1$ and $\sigma_1^{2n+4}\beta_n$ coincide.
    It follows that the operation ${\sigma_1^{2n+4}\beta_n \leadsto \beta_n}$ corresponds to a proper rational tangle replacement $Q_{-2n-4} \leadsto Q_{0}$ from a diagram of $T(3,3n+2)$ to $D_n$. 
\end{proof}

\end{document}
